\originalchapter{介质中的静电场}
\originalsection{极化与束缚电荷}
微观正负电荷在外场中稍微分离, 宏观上用单位体积的电偶极矩 $\vct P$ 表示极化. 对局域偶极分布, 把偶极势 $\vct P(\rr')\cdot(\rr-\rr')/(4\pi\eps|\rr-\rr'|^3)$ 积分并对源坐标分部积分, 得束缚体电荷 $\rho_b=-\nabla\cdot\vct P$ 和表面电荷 $\sigma_b=\vct P\cdot\nn$. 体面两项必须一起计算; 均匀极化球的体电荷为零, 表面电荷却不为零.

令 $\vct D=\eps\ee+\vct P$, 则 $\nabla\cdot\vct D=\rho_f$. $\vct D$ 只把束缚电荷移到构成关系中, 不能独立替代 $\ee$. 对线性各向同性无色散介质, $\vct P=\eps\chi_e\ee$, $\vct D=\varepsilon\ee$, $\varepsilon=\eps(1+\chi_e)$. 这里的模型没有涵盖迟滞、非线性或瞬时各向异性响应.

界面无自由面电荷时 $D_{2n}=D_{1n}$, 静电时切向场连续. 介电常数跳跃会使法向电场跳跃, 不应同时要求 $E_n$ 和 $D_n$ 连续. 边界仍然是电势连续, 但法向导数须按介电常数加权.

\originalsection{介质球与能量}
\begin{example}
介电常数 $\varepsilon$、半径 $a$ 的均匀介质球置于真空均匀场 $E_0\hat{\vct z}$. 求球内场及诱导偶极矩.
\end{example}
\begin{solution}
轴对称且远场为 $\ell=1$, 取 $V_\mathrm{in}=-Ar\cos\theta$, $V_\mathrm{out}=(-E_0r+B/r^2)\cos\theta$. 连续势给 $A=E_0-B/a^3$, 连续 $D_n$ 给 $\varepsilon A=\eps(E_0+2B/a^3)$. 解得
\[
A=\frac{3\eps}{\varepsilon+2\eps}E_0,\qquad
B=a^3\frac{\varepsilon-\eps}{\varepsilon+2\eps}E_0.
\]
内部为均匀场 $A\hat{\vct z}$, 球外偶极矩为 $\pp=4\pi\eps B\hat{\vct z}$. 当 $\varepsilon=\eps$ 时扰动消失, $\varepsilon\to\infty$ 时内场趋零, 对应导体极限.
\end{solution}
对于线性无损静态介质, 缓慢建立自由电荷所需能量为 $U=\frac12\int\ee\cdot\vct D\dd^3x$. 对非线性介质要改成沿建立过程的 $\int\ee\cdot\dd\vct D$; 对色散介质的动态储能也不能直接沿用此式.

介质被电容器吸入的力可由电容随位置的改变计算. 例如忽略边缘场, 面积 $A_d(x)$ 被介质填满、其余面积为真空时,
\[
C(x)=\frac{\eps A+(\varepsilon-\eps)A_d(x)}d.
\]
固定电压时 $F_x=V^2(\varepsilon-\eps)A_d'(x)/(2d)$. 固定电荷时替换 $V=Q/C(x)$ 后仍有趋向增加电容的力, 但力的大小随位置改变. 边缘场负责把局部压力传递成吸入力; 总能量法可以在小间隙近似下避免精细求解边缘.

\originalsection{体源与面源从何而来}
设极化介质占据区域 $\Omega$, $\vct P$ 在内部光滑. 偶极势相加给
$V_b=(4\pi\eps)^{-1}\int_\Omega\vct P(\rr')\cdot\nabla'[1/|\rr-\rr'|]\dd^3x'$.
注意对源坐标求梯度的符号为正, 对观察坐标求梯度则为负. 对 $\nabla'\cdot(\vct P/R)$ 积分,
\[
 V_b=\frac1{4\pi\eps}\left[
 \oint_{\partial\Omega}\frac{\vct P\cdot\nn}{R}\dd a'
 -\int_\Omega\frac{\nabla'\cdot\vct P}{R}\dd^3x'\right].
\]
与普通电荷的库仑势逐项比较, 才能认出 $\rho_b=-\nabla\cdot\vct P$, $\sigma_b=\vct P\cdot\nn$. 把体源积分和面源相加, 总束缚电荷为零, 因为它们来自原来电中性的偶极而非添入净电荷. 若介质只是更大极化系统的一部分, 则须把人工切面的另一侧也计入, 不能把切面电荷误当整系统的新电荷.

以 $\vct P=P_0\hat{\vct z}$ 的均匀极化球为例, 体源为零, 球面 $\sigma_b=P_0\cos\theta$. 表面仅有 $\ell=1$ 模态, 令内势 $Ar\cos\theta$, 外势 $Ba^3\cos\theta/r^2$. 连续势给 $A=B$, 法向场差为 $(2B-(-A))\cos\theta=3A\cos\theta$, 所以 $A=P_0/(3\eps)$. 球内电场 $-\vct P/(3\eps)$, 球外等效偶极矩为体积乘 $\vct P$. 这个负的退极化内场来自两极的束缚电荷, 不表示材料极化率本身为负.

$\vct P$ 是响应, $\vct D$ 是把自由源留在高斯定律中的辅助场. 即便 $\rho_f=0$, 也只能推出 $\nabla\cdot\vct D=0$, 不能推出 $\vct D=0$ 或 $\ee=0$. 比如均匀外场中的介质球, 远处非零场是边界数据, 第5章例题正是无自由体电荷却有非零内部场的情况. 在没有球、柱或平面等足够对称性的几何中, 不能仅从一个高斯曲面的通量决定局部 $\vct D$.


本章依据 8.07 作业6所声明的线性各向同性范围, Quiz 2 的介质题, 并与电容能量推导衔接.
